\chapter{Теория}

\section{Линейная регрессия}
Линейная регрессия (Linear regression) — один из простейших алгоритмов машинного обучения, описывающий зависимость целевой переменной от признака в виде линейной функции:
\[
y = kx + b.
\]

В данном случае была представлена простая (или парная) линейная регрессия, а уравнение вида
\[
f_{w, b}(x) = w_0 x_0 + w_1 x_1 + \dots + w_n x_n + b = \mathbf{w} \cdot \mathbf{x} + b
\]
называется множественной линейной регрессией, где \( b \) — смещение модели, \( \mathbf{w} \) — вектор её весов, а \( \mathbf{x} \) — вектор признаков одного обучающего образца.

К другим условиям определения линейной регрессии относятся гомоскедастичность (дисперсия остатков постоянна и конечна), а также отсутствие мультиколлинеарности (линейной зависимости между признаками).

\begin{figure}[H]
    \centering
    \includegraphics{regr.png}
    \caption{Интуиция линейной регрессии}
\end{figure}

\section{Метод наименьших квадратов и функция потерь}

Выбор регрессионной линии (плоскости), описывающей взаимосвязь данных наилучшим образом, заключается в минимизации функции потерь J, представленной в виде среднеквадратичной ошибки. Проще говоря, линия должна проходить через данные таким образом, чтобы в среднем разница квадратов ожидаемых и реальных значений была минимальна. Данный метод называется методом наименьших квадратов.

При наличии же большого количества выбросов в данных, более эффективным может оказаться метод наименьших модулей, однако у него есть один серьёзный недостаток: функция модуля не является дифференцируемой в точке x=0, что в ряде случаев может затруднить минимизацию ошибки модели. Следовательно, исходя из теоремы Гаусса-Маркова, метод наименьших квадратов является наиболее оптимальной оценкой параметров модели среди всех линейных и несмещённых оценок за счёт меньшей дисперсии.

\section{Датасет}

\begin{longtable}{|p{3cm}|p{2cm}|p{2.5cm}|p{6cm}|p{3cm}|p{3cm}|}
\hline
\textbf{Variable Name} & \textbf{Role} & \textbf{Type} & \textbf{Description} & \textbf{Units} & \textbf{Missing Values} \\
\hline

VendorName & Feature & Categorical & (adviser, amdahl, apollo, basf, bti, burroughs, c.r.d, cambex, cdc, dec, dg, formation, four-phase, gould, honeywell, hp, ibm, ipl, magnuson, microdata, nas, ncr, nixdorf, perkin-elmer, prime, siemens, sperry, stratus, wang) &  & no \\
\hline

ModelName & Feature & Categorical & many unique symbols &  & no \\
\hline

MYCT & Feature & Integer & machine cycle time & nanoseconds & no \\
\hline

MMIN & Feature & Integer & minimum main memory & kilobytes & no \\
\hline

MMAX & Feature & Integer & maximum main memory & kilobytes & no \\
\hline

CACH & Feature & Integer & cache memory & kilobytes & no \\
\hline

CHMIN & Feature & Integer & minimum channels & units & no \\
\hline

CHMAX & Feature & Integer & maximum channels & units & no \\
\hline

PRP & Feature & Integer & published relative performance &  & no \\
\hline

ERP & Feature & Integer & estimated relative performance from the original article &  & no \\
\hline

\end{longtable}


\endinput